% Lösungen Punkte und Strecken im Koordinatensystem · Vierecke nachweisen · Prüfungs-Fokus Abitur GK (zusammenbau v0.6)
\einheitenkopf{Einheit 4 · Vierecke nachweisen}

\erg{1}{a) Parallelogramm und ein rechter Winkel}
\erg{2}{a) $\overrightarrow{AB} = (2 | 2 | 1) = \overrightarrow{DC}$: Parallelogramm; $\overrightarrow{AB} \circ \overrightarrow{AD} = (2 | 2 | 1) \circ (-2 | 1 | 2) = 0$: rechter Winkel, also Rechteck. \quad b) $\overrightarrow{AB} = (3 | 0 | 0) = \overrightarrow{DC}$: Parallelogramm; $\overrightarrow{BC} = (0 | -4 | 5 - p)$, $\overrightarrow{AB} \circ \overrightarrow{BC} = 0$ für jedes $p$: Rechteck; $|\overrightarrow{BC}| = \sqrt{16 + (5 - p)^2}$. \quad c) $\overrightarrow{AB} = \overrightarrow{DC} = (2 | 2 | 1)$: Parallelogramm; $|\overrightarrow{AB}| = |\overrightarrow{AD}| = 3$: alle vier Seiten gleich lang (nicht nur zwei), also Raute. \quad d) $\overrightarrow{AB} = \overrightarrow{DC} = (4 | 0 | 3)$: Parallelogramm; $|\overrightarrow{AB}| = |\overrightarrow{AD}| = 5$: alle vier Seiten gleich lang (nicht nur zwei), also Raute. \quad e) $\overrightarrow{AB} = \overrightarrow{DC} = (1 | 4 | 8)$: Parallelogramm; $|\overrightarrow{AB}| = |\overrightarrow{AD}| = 9$: alle vier Seiten gleich lang (nicht nur zwei), also Raute. \quad f) $\overrightarrow{AB} = (3 | 1 | 0) = \overrightarrow{DC}$: Parallelogramm; $\overrightarrow{AB} \circ \overrightarrow{AD} = 5 \ne 0$: kein rechter Winkel bei $A$, also kein Rechteck. \quad g) $\overrightarrow{AB} = \overrightarrow{DC} = (2 | 1 | 2)$ und $|\overrightarrow{AB}| = |\overrightarrow{AD}| = 3$: Raute; $\overrightarrow{AB} \circ \overrightarrow{AD} = 4 \ne 0$: kein Quadrat. \quad h) $\overrightarrow{AB} = \overrightarrow{DC} = (0 | 3 | 4)$ und $|\overrightarrow{AB}| = |\overrightarrow{AD}| = 5$: Raute; $\overrightarrow{AB} \circ \overrightarrow{AD} = 12 \ne 0$: kein Quadrat. \quad i) $\overrightarrow{EF} = (-6 | 6 | 0) = 2 \cdot \overrightarrow{AB}$: $AB \parallel EF$, also Trapez (nicht aus gleichen Längen schließen); $|\overrightarrow{AE}| = |\overrightarrow{BF}| = \sqrt{18}$. \quad j) $\overrightarrow{PS} = (-6 | 0 | 6) = 2 \cdot \overrightarrow{QR}$: $PS \parallel QR$ und doppelt so lang; $|\overrightarrow{PQ}| = |\overrightarrow{SR}| = \sqrt{93}$.}
\erg{3}{a) $\overrightarrow{AD} = (0 | 6 | 0) = 3 \cdot \overrightarrow{BC}$: Trapez. Höhe ist der Abstand der Parallelen in $x_1$-Richtung, $h = 6$ (nicht $|AB|$). Flächeninhalt $\frac12 \cdot (6 + 2) \cdot 6 = 24$ m². \quad b) $\overrightarrow{EF} = (8 | 0 | 0) = 4 \cdot \overrightarrow{HG}$: Trapez. Höhe $h = 4$ (Abstand in $x_2$-Richtung, nicht der Schenkel $|FG| = 5$). Flächeninhalt $\frac12 \cdot (8 + 2) \cdot 4 = 20$ m². \quad c) $\overrightarrow{KN} = (0 | 6 | 0) = 3 \cdot \overrightarrow{LM}$: Trapez. Höhe $h = 4$ (Abstand in $x_1$-Richtung). Flächeninhalt $\frac12 \cdot (6 + 2) \cdot 4 = 16$ dm². \quad d) $|\overrightarrow{KL}| = |\overrightarrow{KN}| = \sqrt{13}$ und $|\overrightarrow{ML}| = |\overrightarrow{MN}| = 5$: zwei Paare benachbarter gleich langer Seiten (die Gegenseiten sind verschieden lang). \quad e) $|\overrightarrow{KL}| = |\overrightarrow{KN}| = \sqrt{13}$ und $|\overrightarrow{ML}| = |\overrightarrow{MN}| = \sqrt{45}$: zwei Paare benachbarter gleich langer Seiten (die Gegenseiten sind verschieden lang). \quad f) $|\overrightarrow{PQ}| = |\overrightarrow{PS}| = 4$ und $|\overrightarrow{RQ}| = |\overrightarrow{RS}| = \sqrt{10}$: zwei Paare benachbarter gleich langer Seiten (die Gegenseiten sind verschieden lang). \quad g) $\overrightarrow{AE} = (0 | 6 | 1) = 2 \cdot \overrightarrow{CD}$ mit $\overrightarrow{CD} = (0 | 3 | 0{,}5)$: parallel; $\overrightarrow{CD} \circ \overrightarrow{DE} = (0 | 3 | 0{,}5) \circ (-6 | 0 | 0) = 0$: rechter Winkel bei $D$. \quad h) $\overrightarrow{AB} = (2 | 2 | 1) = 2 \cdot \overrightarrow{DC}$ mit $\overrightarrow{DC} = (1 | 1 | 0{,}5)$: parallel; $\overrightarrow{AB} \circ \overrightarrow{AD} = (2 | 2 | 1) \circ (-1 | 2 | -2) = 0$: rechter Winkel bei $A$. \quad i) $\overrightarrow{AB} = (6 | 0 | 8) = 2 \cdot \overrightarrow{DC}$ mit $\overrightarrow{DC} = (3 | 0 | 4)$: parallel; $\overrightarrow{AB} \circ \overrightarrow{AD} = (6 | 0 | 8) \circ (4 | 1 | -3) = 0$: rechter Winkel bei $A$. \quad j) Die Seite an $P$ liegt in der Ebene und steht senkrecht auf $\overrightarrow{PQ} = (0 | 0 | 5)$: Richtung $(4 | -3 | 0)$, ihre Länge $5$ passt schon. $R = P + (4 | -3 | 0) = (6 | -2 | 0)$ (oder $(-2 | 4 | 0)$). \quad k) Die Seite an $P$ liegt in der Ebene und steht senkrecht auf $\overrightarrow{PQ} = (0 | 5 | 0)$: Richtung $(3 | 0 | 4)$, ihre Länge $5$ passt schon. $R = P + (3 | 0 | 4) = (6 | 1 | 8)$ (oder $(0 | 1 | 0)$).}
\erg{4}{a) $B'(3 | b | 0)$ liegt auf der Parallelen zur $x_2$-Achse durch $B$ (nur $x_2$ ändert sich). $(3 | b | 0) = \overrightarrow{AB'}$ und $(3 | b - 7 | 0) = \overrightarrow{CB'}$ sind die beiden Seiten, die in $B'$ zusammenstoßen; Skalarprodukt null heißt, sie stehen senkrecht – rechter Innenwinkel bei $B'$. \quad b) $B'(b | 4 | 0)$ liegt auf der Parallelen zur $x_1$-Achse durch $B$. $(b | 4 | 0) = \overrightarrow{AB'}$ und $(b - 9 | 4 | 0) = \overrightarrow{CB'}$ (Spitze minus Fuß: $B' - C$) sind die Seiten an $B'$; Skalarprodukt null heißt rechter Innenwinkel bei $B'$. \quad c) $P = A + 0{,}36 \cdot \overrightarrow{AB}$ liegt auf der Seite $AB$. Der Inkreismittelpunkt ist der Diagonalenschnittpunkt $M(2 | 3 | 1)$; $\overrightarrow{MP} \circ \overrightarrow{AB} = (1{,}92 | 0 | 1{,}44) \circ (-3 | 0 | 4) = 0$: Der Radius $MP$ steht senkrecht auf der Seite, also berührt der Kreis dort. \quad d) $P = A + 0{,}64 \cdot \overrightarrow{AB}$ liegt auf der Seite $AB$. Der Inkreismittelpunkt ist der Diagonalenschnittpunkt $M(1 | 1 | 2)$; $\overrightarrow{MP} \circ \overrightarrow{AB} = (1{,}92 | 0 | 1{,}44) \circ (3 | 0 | -4) = 0$: Der Radius $MP$ steht senkrecht auf der Seite, also berührt der Kreis dort. \quad e) Die Diagonale $AC$ halbiert die Wand. $P$ liegt zwischen $D$ und $C$, also ist das Dreieck $APD$ kleiner als das Dreieck $ACD$, das die Hälfte ist: Der Teil $ABCP$ ist größer. \quad f) Die Diagonale $BD$ halbiert die Wand. $Q$ liegt zwischen $A$ und $D$, also ist das Dreieck $ABQ$ kleiner als das Dreieck $ABD$, die Hälfte: Der Teil $BCDQ$ ist größer.}
